\chapter{Cubic curves}\label{geometry:ch:cubics}

\begin{center}\itshape
Every curve of the third order is the shadow of one of five parabolas.
\upshape(paraphrase of Newton, \emph{Enumeratio linearum tertii ordinis}, 1704)
\end{center}

The point $(3,5)$ satisfies $y^2=x^3-2$. A tangent through it finds another rational point, $(129/100,383/1000)$. Why should a geometric construction preserve rational coordinates, and how far can we keep going?

Bachet de M\'eziriac described the tangent method around 1621, the year of his Latin Diophantus edition. The arithmetic behind it is simple: after two roots of a cubic polynomial are known, division leaves a linear factor. The geometry makes those two roots visible as a pair of points, or a tangent's double contact. Reflecting the third point leads towards the group law whose complete proof comes in Volume~13, Chapter~15.

Newton's 1704 classification of real cubic curves listed $72$ species and related them to five divergent parabolas. Our normal-form theorem has precise hypotheses: a smooth cubic with a rational flex, over a field of characteristic different from two and three. We will follow its chords, prove the nine-distinct-points associativity case, count flexes and calculate finite-field examples. The prerequisites are Appendices~C and~F and Volume~2's quadratic residues. Full group associativity will remain a clearly identified quoted extension wherever it is mentioned.

\begin{figure}[ht]
\centering
\fig{geometry-ch08-bachet}
\caption{Bachet's tangent on $y^2=x^3-2$. The tangent at $P=(3,5)$, of slope $\frac{27}{10}$, meets the curve again at
$P*P=\bigl(\frac{129}{100},\frac{383}{1000}\bigr)$; reflecting in the $x$-axis gives the point $2P$ of Section~\ref{geometry:sec:weierstrass}.}
\label{geometry:fig:bachet}
\end{figure}

\begin{think}
A line through two rational cubic points leaves a third rational intersection, counted with multiplicity. On a conic there are only two intersections to use. Which part of the construction depends on degree three?
\end{think}

\section{Chords and tangents}

Let $E$ be a smooth plane cubic over a field $K$: its homogeneous equation $F$ has no singular point even over the algebraic closure $\overline K$. It is geometrically irreducible. Indeed a repeated factor makes its component singular, while two distinct components meet by Appendix~F's B\'ezout theorem and their product has zero gradient at an intersection. Thus no line is a component. Write $E(K)$ for the points with coordinates in $K$, and $E(\overline K)$ for the geometric point set.

On a projective line, the restriction of $F$ is a nonzero binary cubic. Its three roots over $\overline K$, counted with their ordinary polynomial root orders, are the line intersections. These one-variable orders are defined directly; they do not require the intrinsic multiplicity theory previewed in Appendix~F.

\begin{lemma}\label{geometry:lem:tangentdouble}
For $P=[p]\in E$ and a line $L$ through it, the restricted cubic has a root of order at least two at $P$ if and only if $L$ is the tangent at $P$.
\end{lemma}
\begin{proof}
Take a second line representative $v$ independent of $p$. The coefficient of $t$ in $F(p+tv)$ is $\nabla F(p)\cdot v$, by expanding monomials; this is valid in every characteristic. The constant coefficient is zero. Thus the root at $t=0$ has order at least two exactly when $v$ belongs to the tangent plane equation $\nabla F(p)\cdot x=0$. That equation is nonzero by smoothness and contains $p$ by Euler's identity $\nabla F(p)\cdot p=3F(p)=0$, also in characteristic three. Its projective line is therefore exactly $L$ in the double-root case.
\end{proof}

For distinct $P,Q\in E(K)$ use their joining line; for $P=Q$ use its tangent. Divide its binary cubic by the two known linear root factors. The remaining factor is linear over $K$, so its root is a point $P*Q\in E(K)$. The operation is commutative. Its defining three-intersection multiset also gives
\[
P*(P*Q)=Q.
\]
This includes repetitions: if $P*Q=P$, the line is tangent at $P$ and its third intersection is $Q$; if all three agree, it is a flex tangent. Hence, for fixed $P$, taking $P*\,$ is an involution, with no excluded tangent cases.

The operation $*$ need not have a neutral element; over the full geometric point set it has none, as an exercise proves. Small finite sets $E(K)$ can be exceptions, so a statement about geometric points must not be transferred automatically to their rational subset. A chosen origin gives a uniform neutral-element construction.

\begin{definition}
Choose $O\in E(K)$ and define $P+Q=O*(P*Q)$ for $P,Q\in E(K)$.
\end{definition}

\begin{proposition}\label{geometry:prop:groupbasic}
This operation is commutative, has neutral element $O$, and gives every $P$ the unique inverse $-P=P*(O*O)$.
\end{proposition}
\begin{proof}
Commutativity follows from $*$. The involution identity gives $P+O=O*(O*P)=P$. Put $O'=O*O$ and $P'=P*O'$. Then $P*P'=O'$, so $P+P'=O*O'=O$. Conversely $P+Q=O$ implies $P*Q=O*O$ by applying $O*\,$, and then $Q=P*(O*O)$ by applying $P*\,$. This proves uniqueness without associativity.
\end{proof}

Translations for $+$ are also bijective: $Q\mapsto P+Q$ is a composite of two involutions. What remains to obtain an abelian group is associativity. The next theorem proves a substantial case with an explicit hypothesis.

\begin{theorem}[associativity for nine distinct points; for all triples it is proved in Volume~13, Chapter~15]\label{geometry:thm:assoc}
For $P,Q,R\in E(K)$ put
\[
S'=P*Q,\quad S=P+Q,\quad T'=S*R,\qquad
U'=Q*R,\quad U=Q+R,\quad T''=P*U.
\]
If $P,Q,R,O,S,S',U,U',T'$ are nine distinct points, then $(P+Q)+R=P+(Q+R)$.
\end{theorem}
\begin{proof}
The two sides are $O*T'$ and $O*T''$, so it suffices to prove $T'=T''$. Consider the two line-union cubics
\[
\begin{aligned}
C_1&=(PQS')(SRT')(U'OU),\\
C_2&=(S'OS)(QRU')(PUT''),
\end{aligned}
\]
where a parenthesized triple names its common line. Each line of $C_1$ contains exactly its three displayed points of $E$; the distinctness hypothesis makes the nine points different and the three lines different. Thus $C_1$ and the irreducible $E$ have no common component and meet in those nine distinct points. The cubic $C_2$ contains eight of them, so Cayley--Bacharach forces it to contain $T'$. Its first line has only $S',O,S$ on $E$, and its second only $Q,R,U'$, so neither contains $T'$. Therefore $T'$ is on its third line $PU$, whose third intersection is $T''$. Since $T'$ differs from $P,U$, it equals $T''$, proving the claim.
\end{proof}

\begin{figure}[ht]
\centering
\fig{geometry-ch08-real}
\caption{Two real cubics: $y^2=x^3-x$ has two components, $y^2=x^3-x+1$ has one. The rust points are the real inflection points in the affine plane; with
the point at infinity $O$ they are the three points of the vertical dashed line (Theorem~\ref{geometry:thm:realflexes}).}
\label{geometry:fig:realcubics}
\end{figure}

\begin{remark}
Coincident intersections, including the tangent cases, are not covered by that theorem. We quote the full group-law result from Volume~13, Chapter~15, where the intersection theorem handles multiplicities. All computations proved below use the chord identities directly, retain parentheses, or verify their own finite table; they do not infer unrestricted associativity from the nine-distinct case.
\end{remark}

\section{The Weierstrass form}\label{geometry:sec:weierstrass}

A point $O$ is a \emph{flex}, or inflection point, when its tangent has threefold contact there; equivalently $O*O=O$. With a flex as origin the inverse becomes $-P=P*O$. The collinearity exercise will also show directly that $(P+Q)+R=O$ precisely when $P,Q,R$ are the three intersections of a line, counted with multiplicity.

\begin{proposition}[Weierstrass form]\label{geometry:prop:weierstrass}
Suppose $\operatorname{char}K\ne2,3$ and a smooth cubic has a flex $O\in E(K)$. There are projective coordinates over $K$ in which $O=[0:1:0]$ and
\[
y^2z=x^3+axz^2+bz^3,\qquad 4a^3+27b^2\ne0.
\]
Conversely each such equation defines a smooth cubic with that flex. In the affine chart it is $y^2=x^3+ax+b$.
\end{proposition}
\begin{proof}
Choose a basis adapted to $O$ and its tangent, making them $[0:1:0]$ and $z=0$. The restriction to that tangent has a triple root at $O$, so $F(x,y,0)=cx^3$ with $c\ne0$, since a tangent is not a component. Hence $F=cx^3+zQ$. The $y^2$ coefficient $e$ in $Q$ is nonzero: otherwise every first partial would vanish at $O$. Write
\[
Q=ey^2+y(\alpha x+\beta z)+q(x,z).
\]
Set $y_{\rm old}=y_{\rm new}-(\alpha x+\beta z)/(2e)$ to complete the square. Rearranging gives
\[
ey^2z=c'x^3+c_2x^2z+c_1xz^2+c_0z^3,\quad c'=-c\ne0.
\]
Substitute $z_{\rm old}=(c'/e)z_{\rm new}$ and divide by $c'$. Both leading coefficients become one. If the new $x^2z$ coefficient is $s$, the substitution $x_{\rm old}=x_{\rm new}-sz/3$ removes it. These invertible changes fix $O$ and give the asserted short form.

Let $f=x^3+ax+b$. At an affine singularity, $2y=0$ and $f'=0$, with $y^2=f$. Thus $y=0$ and $f,f'$ have a common root. Such a root satisfies $x^2=-a/3$ and $2ax/3+b=0$. If $a=0$, this means $b=0$; otherwise it means $x=-3b/(2a)$ and $27b^2=-4a^3$. Conversely these equalities give a common root, including $a=b=0$. Therefore smoothness is equivalent to $4a^3+27b^2\ne0$. The only infinity point is $O$, where $F_z=1$, so it is always smooth. Restricting to $z=0$ gives $-x^3$, proving its flex property and the converse.
\end{proof}

On this curve $O$ joins $(x,y)$ by its vertical line, whose other intersection is $(x,-y)$. Hence negation is reflection in the horizontal axis. To add two affine points, take their third line intersection and reflect it. If $Q=-P$ the answer is $O$. Otherwise the required slope is
\[
\lambda=\frac{y_2-y_1}{x_2-x_1}\quad(x_1\ne x_2),\qquad
\lambda=\frac{3x_1^2+a}{2y_1}\quad(P=Q).
\]
In the second case $y_1\ne0$, since the excluded case is self-inverse. Substituting the line into the cubic and comparing its $x^2$ coefficient gives
\[
x_3=\lambda^2-x_1-x_2,\qquad
P+Q=(x_3,-y_1-\lambda(x_3-x_1)).
\]
Sums involving $O$ use its neutral-element property. Thus vertical, tangent and infinity cases all have a rule.

For Bachet's point, $\lambda=27/10$ and
\[
2P=\left(\frac{129}{100},-\frac{383}{1000}\right).
\]
The third tangent intersection before reflection has positive ordinate $383/1000$. Its exact check is $129^3-2\cdot100^3=146689=383^2$.

Repeated doubling already gives infinitely many rational points here, without full associativity. For a nonzero rational $r$, let $v_2(r)$ be the exponent of two in its prime factorization. Products add exponents; a sum with one strictly least exponent retains it, by factoring out that power of two and observing the remaining numerator is odd. The doubling abscissa on this curve is
\[
x_D=\frac{x^4+16x}{4(x^3-2)}.
\]
If $v_2(x)<0$, the numerator has exponent $4v_2(x)$ and the denominator $2+3v_2(x)$, so $v_2(x_D)=v_2(x)-2$. Starting from $x(2P)=129/100$, of exponent $-2$, these exponents decrease forever. Also $x^3-2\ne0$ at every step by its negative exponent, so the tangent denominator is nonzero. The resulting rational points are all distinct. The denominator exercise proves the same argument for any integral short equation.

\begin{remark}
The full group law gives a second, deeper interpretation. We quote Nagell--Lutz: on a nonsingular integral short equation a rational torsion point has integer coordinates. It would imply Bachet's point has infinite order because its double is nonintegral. We also quote Mordell's 1922 finite-generation theorem for $E(\mathbb Q)$, following Poincar\'e's 1901 question, and Mazur's 1977 list of possible torsion groups: $\mathbb Z/n$ for $1\le n\le10$ or $n=12$, and $\mathbb Z/2\times\mathbb Z/2m$ for $1\le m\le4$. Volume~13 proves the first two results; Mazur's theorem lies beyond this series' elementary arguments. The free abelian rank is central to the Birch--Swinnerton-Dyer conjecture. None of these quoted results is needed for the rational constructions just proved.
\end{remark}

\section{Nine flexes and twelve lines}

\begin{think}
At a flex, a tangent's third intersection returns to its contact point. On the real curves in Figure~\ref{geometry:fig:realcubics}, can the two affine flexes share an $x$-coordinate? Where does their vertical line meet the projective cubic a third time?
\end{think}

The flex condition is $P*P=P$. With the quoted abelian group law this is also $3P=O$; the following count uses only the tangent formula, not that group interpretation.

\begin{theorem}\label{geometry:thm:nineflexes}
A smooth short Weierstrass cubic over an algebraically closed field of characteristic zero has exactly nine flexes.
\end{theorem}
\begin{proof}
Put $f=x^3+ax+b$. A point $(x,0)$ has a vertical tangent with third intersection $O$, so is not a flex. For $y\ne0$, the tangent slope is $f'(x)/(2y)$ and its third intersection has abscissa $\lambda^2-2x$. Because the tangent is nonvertical, that intersection equals $P$ precisely when its abscissa is $x$. The condition is therefore
\[
\psi(x)=12xf(x)-f'(x)^2
=3x^4+6ax^2+12bx-a^2=0.
\]
Differentiation gives $\psi'=12f$, since $f''=6x$. A multiple root of $\psi$ would have $f=0$ and then $f'=0$, contradicting smoothness. No root of $\psi$ is a root of $f$ for the same reason. Thus the quartic has four distinct roots, each giving two distinct nonzero ordinates. Those eight affine flexes together with $O$ give exactly nine.
\end{proof}

\begin{remark}
Under the full abelian group theorem quoted from Volume~13, these points form $E[3]$, a nine-element group annihilated by three, hence isomorphic to $(\mathbb Z/3)^2$. The third intersection of a line through distinct flexes is then $-P-Q$, another flex. The twelve three-point lines are the affine lines of that two-dimensional vector space over $\mathbb F_3$, counted by $\binom92/\binom32=12$. This group argument is quoted with its prerequisite; it is not an additional proof of full associativity. The Hesse-pencil exercise below verifies the explicit nine-point configuration independently through equations.
\end{remark}

This is the \emph{Hesse configuration}, with nine points and twelve lines. It is the affine plane of order three, obtained by removing a line and its four points from $\mathbb P^2(\mathbb F_3)$, not the thirteen-point projective plane itself. Maclaurin knew the nine flexes; Hesse studied their configuration. Artebani and Dolgachev give the historical account. A real straight-line drawing of the whole configuration is impossible by Appendix~C's Sylvester--Gallai theorem: every determined line would contain three points. The existing drawing therefore uses circles for four of its combinatorial lines.

\begin{figure}[ht]
\centering
\fig{geometry-ch08-hesse}
\caption{The Hesse configuration: nine points, twelve lines of three points each, any two points on exactly one line. Eight lines are straight; the four
other ones are drawn as circles.}
\label{geometry:fig:hesse}
\end{figure}

\begin{theorem}\label{geometry:thm:realflexes}
A smooth real short Weierstrass cubic has exactly three real flexes, and they are collinear.
\end{theorem}
\begin{proof}
Again $\psi'=12f$, and at a real root $e$ of the squarefree cubic $f$ we have $\psi(e)=-f'(e)^2<0$. Also $\psi(x)\to+\infty$ at both ends of the real line.

If $f$ has one real root $e$, the derivative $\psi'$ is negative before it and positive after it. Thus $\psi$ has exactly two real roots, one on each side of $e$. The left root has $f<0$ and gives no real cubic point; the right root has $f>0$ and gives the two real affine flexes with opposite ordinates.

If $f$ has three real roots $e_1<e_2<e_3$, the derivative signs are minus, plus, minus, plus on the four intervals. The values at all three roots are negative. Consequently $\psi$ has one root before $e_1$, none between the roots, and one after $e_3$. Again only the rightmost root has $f>0$, giving two real affine flexes. In both cases the point $O$ supplies the third real flex, and all three lie on the vertical line through those two affine points.
\end{proof}

\begin{remark}
We quote the extension to every smooth complex plane cubic: it has nine flexes, found as its nine simple intersections with the Hessian cubic, the determinant of the matrix of second partial derivatives. In particular a complex flex exists, allowing the Weierstrass normal-form theorem to apply. Kirwan and Brieskorn--Kn\"orrer prove this Hessian statement. The short-form count above did not assume its proof.
\end{remark}

\section{Cubics over finite fields}

Let $p$ be an odd prime and $f$ a monic squarefree cubic over $\mathbb F_p$. The smooth curve $y^2=f(x)$ has the single infinity point $O$. This includes the short equations considered above whenever they are smooth, even if $p=3$. For $c\in\mathbb F_p$, the Legendre symbol $\left(\frac cp\right)$ is one for a nonzero square, minus one for a nonsquare, and zero for zero.

\begin{proposition}\label{geometry:prop:countEp}
For this curve,
\[
\#E(\mathbb F_p)=p+1+\sum_{x\in\mathbb F_p}\left(\frac{f(x)}p\right).
\]
\end{proposition}
\begin{proof}
For each fixed $x$, the number of solutions of $y^2=f(x)$ is $1+\left(\frac{f(x)}p\right)$: respectively two, zero or one. Summing gives the affine count, and adding the unique point $O$ gives the formula.
\end{proof}

The sum consists of correlated residue values; it is not a sequence of independent coin flips. We quote Hasse's 1933 theorem that its absolute value is at most $2\sqrt p$, proved in Volume~13, Chapter~16. The displayed curve over $\mathbb F_{101}$ has $99$ affine points, hence $100$ in all, within that bound. Some residue sums cancel exactly by a simple symmetry.

\begin{proposition}\label{geometry:prop:supersingular}
If $p\equiv3\pmod4$, the curve $y^2=x^3-x$ has $p+1$ points over $\mathbb F_p$.
\end{proposition}
\begin{proof}
Its cubic is squarefree for every odd $p$: its roots $0,1,-1$ are distinct. Since $f(-x)=-f(x)$ and $-1$ is a nonsquare for $p\equiv3\pmod4$, the Legendre values at $x,-x$ are opposite, including zero values. All nonzero $x$ pair off and the value at zero is zero. Thus the residue sum vanishes and the preceding formula gives $p+1$.
\end{proof}

\begin{figure}[ht]
\centering
\fig{geometry-ch08-finite}
\caption{The $99$ affine points of $y^2=x^3-x+1$ over $\mathbb F_{101}$; with $O$, $E(\mathbb F_{101})$ has $100$ points, and $|100-102|=2\le2\sqrt{101}$.}
\label{geometry:fig:finitecubic}
\end{figure}

\begin{remark}
Koblitz and Miller independently proposed elliptic-curve cryptography in 1985. Using the quoted group theorem, the chord formulas give efficient addition in suitable finite groups, whereas recovering an integer multiplier from a point and its multiple can be difficult for appropriately chosen parameters. This does not say every curve or small example is secure. The examples here are small enough to check by hand; their purpose is to make the arithmetic visible.
\end{remark}

\FloatBarrier
\begin{figure}[ht]
\centering
\fig{v9-appg-modular-chord}
\caption{Addition on $y^2=x^3+x+1$ over $\mathbb F_5$. The modular line $y=4x+1$ contains $P=(0,1)$, $Q=(4,2)$ and $R=(2,4)$ on the curve. Reflecting $R$ modulo five gives $P+Q=(2,1)$. The broken drawing reflects reduction modulo five; the line is not one ordinary Euclidean segment in this grid.}
\label{fig:v9-modular-chord}
\end{figure}
\begin{figure}[ht]
\centering
\fig{v9-appg-hesse-pencil}
\caption{The Hesse pencil $X^3+Y^3+Z^3-3\kappa XYZ=0$. The three finite parameters $1,\omega,\omega^2$ with $\omega^3=1$ give triangles, as does the member $XYZ=0$ at infinity. Every other member is smooth; all share the same nine coordinate-line basepoints, which are its flexes.}
\label{fig:v9-hesse-pencil}
\end{figure}
\FloatBarrier

A cubic's third intersection has given us rational constructions, a partially proved addition law and exact point counts. The flex quartic separates its real and complex geometry, while finite fields turn the formulas into small arithmetic tables. The exercises finish several computations without borrowing full associativity. Appendix~H then returns to Euclidean space, where integer combinations of vectors form lattices and geometric questions become arithmetic ones.
\FloatBarrier

\section{Exercises}

\begin{exercise}[$\star\star$]\label{geometry:ex:noneutral}
Prove that $*$ has no neutral element on the geometric point set of a smooth cubic over an algebraically closed field. Distinguish this from $E(K)$: show that the smooth cubic
\[
y^2z+yz^2=x^3+xz^2+z^3
\]
over $\mathbb F_2$ has only $O=[0:1:0]$ as a rational point, and that $O$ is neutral for $*$ on that singleton set.
\end{exercise}
\begin{exercise}[$\star$]\label{geometry:ex:collinear}
With a flex as origin, prove $-P=P*O$ and
$(P+Q)+R=O$ exactly when $P,Q,R$ are a line's three intersections with the cubic, counted with multiplicity. Use no associativity assumption.
\end{exercise}
\begin{exercise}[$\star$]\label{geometry:ex:bachet}
For $P=(3,5)$ on $y^2=x^3-2$, compute $3P=(P+P)+P$. Verify the result by an explicit integer identity after clearing denominators.
\end{exercise}
\begin{exercise}[$\star$]\label{geometry:ex:twotorsion}
In characteristic different from two, let $y^2=f(x)$ with $f$ a monic squarefree cubic and $O$ its infinity point. Prove that the affine solutions of $P+P=O$ are exactly $(e,0)$ with $f(e)=0$. Over $\mathbb C$, verify directly that these three points and $O$ have the addition table of $(\mathbb Z/2)^2$. For real coefficients, when are all four real?
\end{exercise}
\begin{exercise}[$\star\star$]\label{geometry:ex:torsionsix}
Over $\mathbb Q$, on the curve $y^2=x^3+1$, start at $P=(2,3)$ and define $1P=P$, $nP=(n-1)P+P$. Compute the first six terms and prove the first return to $O$ is at $6P$. Interpret this as order six under the quoted full group law.
\end{exercise}
\begin{exercise}[$\star$]\label{geometry:ex:realcomponents}
For a smooth real short cubic $y^2=f(x)$, prove its projective real point set has two connected components if $f$ has three real roots and one if it has one. Account for the point at infinity and its local arc.
\end{exercise}
\begin{exercise}[$\star\star$]\label{geometry:ex:fermatflexes}
Find all nine flexes of the complex Fermat cubic $x^3+y^3+z^3=0$. Prove directly that the line through any two distinct flexes contains a third flex.
\end{exercise}
\begin{exercise}[$\star\star$]\label{geometry:ex:flexesxcube}
Determine all complex flexes of $y^2=x^3+1$, identify its real flexes and check their collinearity. Verify the three-step return of $(0,1)$ directly from tangents and chords.
\end{exercise}
\begin{exercise}[$\star$]\label{geometry:ex:countfive}
List $E(\mathbb F_5)$ for $y^2=x^3+x+1$. For $P=(0,1)$ and recursively parenthesized multiples $nP=(n-1)P+P$, show the first nine terms visit every point and $9P=O$. Determine all rational flexes. State separately the cyclic-group interpretation that uses the full group theorem.
\end{exercise}
\begin{exercise}[$\star\star$]\label{geometry:ex:xcubeplusone}
If the prime $p\equiv2\pmod3$, prove the projective equation $y^2z=x^3+z^3$ has $p+1$ rational points. Include $p=2$ and explain whether that case is a smooth elliptic curve.
\end{exercise}
\begin{exercise}[$\star\star$]\label{geometry:ex:changeorigin}
Let $O$ be a flex and $O_1$ any rational cubic point. Write $+$ for the operation with origin $O$ and $\oplus$ for that with origin $O_1$. Prove, with parentheses retained,
\[
X*Y=-(X+Y),\qquad P\oplus Q=-\bigl(O_1+(-(P+Q))\bigr).
\]
Explain how the quoted full group theorem turns this into $P\oplus Q=P+Q-O_1$ and gives an isomorphism from the new-origin group to the old one.
\end{exercise}
\begin{exercise}[$\star\star\star$]\label{geometry:ex:sixonconic}
Let $O$ be a flex and six distinct cubic points $P_1,\ldots,P_6$ lie on a conic $C$. Put $A=P_1*P_2$, $B=P_3*P_4$, $C'=P_5*P_6$. Assume these nine points are distinct and $C'\notin C$. Prove $A,B,C'$ collinear and $(A+B)+C'=O$. Identify what extra theorem is needed to express this as the unparenthesized six-point sum being zero.
\end{exercise}

\begin{exercise}[$\star\star\star$]\label{ex:v9-hesse-pencil}
For $\kappa\in\mathbb C$, study $E_\kappa:X^3+Y^3+Z^3-3\kappa XYZ=0$. Prove it is smooth exactly when $\kappa^3\ne1$. Factor the three singular members and the member $XYZ=0$ at parameter infinity. Show all smooth members have the same nine flexes on the coordinate lines. Prove the twelve lines in those four triangles give the nine-point affine-plane incidence pattern, without using full group associativity.
\end{exercise}
\begin{exercise}[$\star\star$]\label{ex:v9-quadratic-twist-count}
Let $p$ be odd, $f$ a monic squarefree cubic and $d$ a nonsquare in $\mathbb F_p$. Compare $E:y^2=f(x)$ and its \emph{quadratic twist} $E_d:dy^2=f(x)$. Prove $\#E+\#E_d=2(p+1)$. For $p=5$, $f=x^3+x+1$ and $d=2$, list the twist's points and compare its count with nine. Give an isomorphism after adjoining $\sqrt d$ and explain why it cannot be defined over $\mathbb F_5$ in this example.
\end{exercise}
\begin{exercise}[$\star\star$]\label{ex:v9-denominator-doubling}
On a nonsingular rational short equation $y^2=x^3+ax+b$ with $a,b\in\mathbb Z$, let $D(P)=P+P$. Derive
\[
x(D(P))=\frac{x^4-2ax^2-8bx+a^2}{4(x^3+ax+b)}.
\]
For rational $P$ with $v_2(x)<0$, prove $v_2(x)$ is even, $v_2(y)=3v_2(x)/2$, and $v_2(x(D(P)))=v_2(x)-2$. Deduce that all repeated doublings are defined, rational and distinct without using associativity or Nagell--Lutz. Apply this to $2(3,5)$ on Bachet's curve and give both coordinate valuations after $n$ further doublings.
\end{exercise}
\begin{exercise}[$\star\star\star$]\label{ex:v9-nodal-multiplication}
Over a field of characteristic different from two, remove the node $S=[0:0:1]$ from $N:y^2z=x^2(x+z)$, and use the smooth flex $O=[0:1:0]$ as origin. Prove the chord operation stays on the smooth points and the addition law is isomorphic to multiplication in $K^\times$. Use the parameter $x=t^2-1$, $y=t(t^2-1)$ and the new coordinate $u=(t+1)/(t-1)$, treating the projective parameter and vertical/tangent cases explicitly.
\end{exercise}

\section{Hints}

\begin{hint}{geometry:ex:noneutral}
A proposed neutral $N$ is a flex. Adapt coordinates to it and its tangent, then show its polar quadratic is nonzero, including characteristic two. Infinitely many tangents through it would contradict B\'ezout. For the finite example compare $y^2+y$ with $x^3+x+1$ over $\mathbb F_2$.
\end{hint}
\begin{hint}{geometry:ex:collinear}
Apply the involution $O*\,$ and then the involution $(P+Q)*\,$ to the equation; keep the original parentheses.
\end{hint}
\begin{hint}{geometry:ex:bachet}
Use $2P=(129/100,-383/1000)$. The next chord's slope is $5383/1710$.
\end{hint}
\begin{hint}{geometry:ex:twotorsion}
The inverse is $(x,-y)$ and is unique. Add two distinct root points along $y=0$ to determine the whole four-point table.
\end{hint}
\begin{hint}{geometry:ex:torsionsix}
Compute the tangent at $(2,3)$, then repeatedly add $(2,3)$. The last chord is vertical.
\end{hint}
\begin{hint}{geometry:ex:realcomponents}
Use the intervals where $f\ge0$. Near $O$ the chart $y=1$ has coordinates $u=x/y,v=z/y$ and equation $v=u^3+auv^2+bv^3$.
\end{hint}
\begin{hint}{geometry:ex:fermatflexes}
The coordinate lines supply three points each. Verify triple tangency at one and use coordinate symmetries. For the cross-family lines use a determinant and cube-root exponents modulo three.
\end{hint}
\begin{hint}{geometry:ex:flexesxcube}
The flex polynomial is $3x(x^3+4)$. At $(0,1)$ the horizontal tangent's third contact is itself.
\end{hint}
\begin{hint}{geometry:ex:countfive}
The squares are $0,1,4$. In characteristic five the flex polynomial is $3x^4+x^2+2x+4$; test every abscissa actually present on the curve.
\end{hint}
\begin{hint}{geometry:ex:xcubeplusone}
Cubing is bijective when $\gcd(3,p-1)=1$. Count by choosing $y$ first; the argument also works at $p=2$.
\end{hint}
\begin{hint}{geometry:ex:changeorigin}
Negation is $X\mapsto X*O$. Apply it to the defining sum, then rewrite the new-origin definition once more.
\end{hint}
\begin{hint}{geometry:ex:sixonconic}
Compare $E$ with $P_1P_2\cup P_3P_4\cup P_5P_6$ and use $C\cup AB$ as the third cubic.
\end{hint}
\begin{hint}{ex:v9-hesse-pencil}
Multiply the three zero-gradient equations when all coordinates are nonzero. On the tangent at $[1:-u:0]$, $u^3=1$, substitute $[1:-u-\kappa u^2t:t]$.
\end{hint}
\begin{hint}{ex:v9-quadratic-twist-count}
The two affine counts for each $x$ are $1+\chi(f(x))$ and $1-\chi(f(x))$. Both curves have one point at infinity.
\end{hint}
\begin{hint}{ex:v9-denominator-doubling}
Use the slope formula to combine the doubling abscissa into one fraction. In each numerator or denominator, show the leading $x$-power has strictly smaller valuation than every other term.
\end{hint}
\begin{hint}{ex:v9-nodal-multiplication}
A nonvertical chord gives the parameter polynomial $t^3-mt^2-t+m-b$. Compare its values at $1,-1$. A chord with two smooth intersections cannot also pass through the double node.
\end{hint}

\section{Solutions}

\begin{solution}{geometry:ex:noneutral}
\idea A geometric neutral would force infinitely many contacts on its nonzero polar; a rational singleton does not have infinitely many points to test.
\writeup Suppose $N*P=P$ for every geometric point $P$. Taking $P=N$ makes $N$ a flex. Choose coordinates $N=[0:1:0]$ with tangent $z=0$. As in the beginning of the normal-form proof, without dividing by two or three,
\[
F=cx^3+z\bigl(ey^2+y(\alpha x+\beta z)+q_2x^2+q_1xz+q_0z^2\bigr),
\]
with $c,e\ne0$. Its polar for $N$ is $F_y=z(2ey+\alpha x+\beta z)$. Outside characteristic two it is nonzero. In characteristic two, if it were zero then $\alpha=\beta=0$. Set $z=1$, choose $x$ with $cx^2+q_1=0$, and choose $y$ with $ey^2+q_2x^2+q_0=0$. Algebraic closure supplies both roots. These equations make $F_x=F_y=F_z=0$, and Euler's identity, with $3=1$ in this characteristic, also makes $F=0$. This contradicts smoothness. Thus the polar is nonzero in every characteristic.

For every $P\ne N$, the equation $N*P=P$ makes $NP$ tangent at $P$, so $F_y(P)=0$. The smooth cubic has infinitely many geometric points by Appendix~F, while its nonzero degree-two polar has no common component with it. B\'ezout allows at most six intersections, contradiction. Hence no geometric neutral exists.

For the binary example, bring all terms to one side. Its partials are $x^2+z^2,z^2,y^2+z^2$, with no simultaneous projective zero even over the algebraic closure, so the curve is smooth. Affinely $y^2+y=0$ for either binary $y$, but $x^3+x+1=1$ for either binary $x$, so there are no affine rational points. At infinity $x^3=0$ gives only $O$. Its tangent $z=0$ has triple contact there, so $O*O=O$. Thus $O$ is neutral on $E(\mathbb F_2)=\{O\}$, explicitly disproving the unqualified rational-point claim.
\end{solution}
\begin{solution}{geometry:ex:collinear}
\idea The involution identities remove the need to reassociate anything.
\writeup The inverse formula becomes $-P=P*(O*O)=P*O$. If $P,Q,R$ are the line intersections, then $P*Q=R$ and $P+Q=O*R=-R$. Thus $(P+Q)+R=(-R)+R=O$. Conversely from $O*((P+Q)*R)=O$, applying $O*\,$ gives $(P+Q)*R=O*O=O$. Applying $(P+Q)*\,$ gives $R=(P+Q)*O$. The definition $P+Q=O*(P*Q)$ and the involution identity give $(P+Q)*O=P*Q$. Hence $R=P*Q$, which is exactly the line-intersection condition, including tangent and flex repetitions. No parentheses have been changed.
\end{solution}
\begin{solution}{geometry:ex:bachet}
\idea One exact chord calculation gives the next rational point, and a cleared integer equality verifies it.
\writeup The slope from $P$ to $2P$ is $(-5383/1000)/(-171/100)=5383/1710$. Substituting into the addition formula gives
\[
3P=\left(\frac{164323}{29241},-\frac{66234835}{5000211}\right),
\qquad29241=171^2,\quad5000211=171^3.
\]
The curve check is the exact identity
\[
66234835^2=164323^3-2\cdot171^6
=4387053367477225.
\]
Dividing by $171^6$ gives $y^2=x^3-2$. This proves the requested substitution directly.
\end{solution}
\begin{solution}{geometry:ex:twotorsion}
\idea Unique inverses identify the points, and their finite table proves its own group structure.
\writeup Since the inverse is $(x,-y)$ and is unique, $P+P=O$ exactly when $P=-P$. For an affine point in characteristic different from two this is $y=0$, with $f(x)=0$. Over $\mathbb C$ squarefreeness gives three distinct roots $e_1,e_2,e_3$. Put $T_i=(e_i,0)$. We have $T_i+T_i=O$; for distinct $i,j$, the line $y=0$ has third intersection $T_k$ and $O*T_k=T_k$, so $T_i+T_j=T_k$. Along with neutral $O$, this is precisely addition on $\{(0,0),(1,0),(0,1),(1,1)\}$ over $\mathbb F_2$, with $T_1,T_2,T_3$ assigned to the three nonzero vectors. It therefore proves an associative four-point subgroup directly, independently of the larger operation. With real coefficients, all four are real exactly when all three cubic roots are real.
\end{solution}
\begin{solution}{geometry:ex:torsionsix}
\idea Follow the prescribed recursion, so the computation proves exactly the stated first return.
\writeup At $(2,3)$ the tangent slope is two, giving $2P=(0,1)$. The chord slope from that point to $P$ is one, giving $3P=(-1,0)$. The next slopes are one and two, giving $4P=(0,-1)$ and $5P=(2,-3)$. Finally the vertical chord from $5P$ to $P$ has third intersection $O$, so $6P=O*O=O$. Each displayed point is on the curve by substitution, and none of the first five is $O$. These are the recursively defined terms without any associativity assumption; under the quoted group theorem the same list proves order six.
\end{solution}
\begin{solution}{geometry:ex:realcomponents}
\idea Join the upper and lower graphs at their finite roots and inspect their common infinity point locally.
\writeup If $f$ has three simple real roots $e_1<e_2<e_3$, then $f\ge0$ on $[e_1,e_2]\cup[e_3,\infty)$. The two square-root graphs on the bounded interval join at both ends into one connected closed oval. On the unbounded interval they join at $e_3$ and both approach $O$, since $x/y\to0$ and $1/y\to0$ as $x\to\infty$. Near $O$, the equation in the chart $y=1$ is $v=u^3+auv^2+bv^3$. Its derivative with respect to $v$ at zero is one after moving terms to one side, so the implicit-function theorem gives one smooth arc $v=v(u)$. Its two sides join precisely the upper and lower unbounded ends. Thus that projective branch is connected and is the second component. The bounded oval stays separate: its compact $x$-interval is disjoint from $[e_3,\infty)$ and it has no limit at $O$. These two compact disjoint pieces exhaust the real set. If $f$ has only one real root, only the unbounded interval occurs, and the same local argument gives exactly one component.
\end{solution}
\begin{solution}{geometry:ex:fermatflexes}
\idea Coordinate symmetries find the flexes; their exponent relations find every three-point line.
\writeup Let $\omega$ be a primitive cube root of unity. The nine coordinate-line points are
\[
[1:-\omega^i:0],\quad[0:1:-\omega^j],\quad[-\omega^k:0:1],
\quad i,j,k=0,1,2.
\]
The cubic is smooth because its gradient $(3x^2,3y^2,3z^2)$ never vanishes projectively. At $[1:-1:0]$ the tangent is $x+y=0$, and restriction of the cubic to it is $z^3$, proving triple contact. Coordinate permutations and multiplication by cube roots of unity carry this point to all nine and preserve the equation. Once one flex is known, the normal-form theorem applies, and the nine-flex theorem proves these are all the flexes.

Two in one family have their third on the same coordinate line. For different families the determinant of the three displayed representatives is $1-\omega^{i+j+k}$. Given two exponents from different families, the unique third exponent with $i+j+k=0$ modulo three makes them collinear. Thus every pair of distinct flexes has a third flex on its line, by direct coordinates rather than group associativity.
\end{solution}
\begin{solution}{geometry:ex:flexesxcube}
\idea The quartic separates the two real affine flexes from the six nonreal ones.
\writeup Here $\psi=3x(x^3+4)$. The root zero gives $(0,1),(0,-1)$. Write $c=\sqrt[3]4>0$. The other roots $-c,-c\omega,-c\omega^2$ each give $y^2=-3$, hence the six points with ordinates $\pm i\sqrt3$. With $O$ these are nine distinct flexes. Only $O,(0,1),(0,-1)$ are real, all on $x=0$. For $Q=(0,1)$, the horizontal tangent $y=1$ has triple intersection at $Q$, so $Q*Q=Q$ and $2Q=O*Q=(0,-1)$. The vertical chord from $2Q$ to $Q$ gives $3Q=O$. This verifies the return without reassociating $2P$ from the earlier exercise.
\end{solution}
\begin{solution}{geometry:ex:countfive}
\idea Compute the finite list, then test the flex equation at every allowable abscissa.
\writeup The cubic values at $0,1,2,3,4$ are $1,3,1,1,4$ modulo five. Thus the points are $O$, $(0,\pm1),(2,\pm1),(3,\pm1),(4,\pm2)$. The chord formulas, always reducing modulo five, give the successive list
\[
\begin{gathered}
(0,1),\ (4,2),\ (2,1),\ (3,-1),\ (3,1),\\
(2,-1),\ (4,-2),\ (0,-1),\ O.
\end{gathered}
\]
For example the first two slopes are $1/2=3$ and $1/4=4$, giving $(4,2)$ and $(2,1)$; the remaining chord slopes are respectively $0,1,0,4,3$ before the final vertical chord. Substituting each into the formula verifies the list. Its points are distinct and exhaust all nine, with first return at nine. Figure~\ref{fig:v9-modular-chord} shows its second addition.

For affine points present here $y\ne0$, the tangent flex condition is $\psi=3x^4+x^2+2x+4=0$. At abscissae $0,2,3,4$ its values are $4,0,2,1$. Thus the rational flexes are precisely $O,(2,1),(2,-1)$. This checks every affine point, without inferring multiples from associativity. Once the full group law is supplied, the nine-term orbit says the group is cyclic of order nine.
\end{solution}
\begin{solution}{geometry:ex:xcubeplusone}
\idea Choose an ordinate and solve uniquely for its cubed abscissa.
\writeup Since $3$ is coprime to $p-1$, cubing is bijective on $\mathbb F_p^\times$, by its cyclic group structure, and also on the whole field after including zero. For each of the $p$ choices of $y$, there is therefore exactly one $x$ with $x^3=y^2-1$. This gives $p$ affine points, plus the single infinity point $O$. The argument includes $p=2$, where the points are $(0,1),(1,0),O$. In that characteristic the equation is singular at $(0,1)$: the affine partials of $y^2-x^3-1$ are $x^2,0$, both zero there. The point count is still valid, but this case is not a smooth elliptic curve.
\end{solution}
\begin{solution}{geometry:ex:changeorigin}
\idea Rewrite each chord through the involution given by the old flex origin.
\writeup From $X+Y=O*(X*Y)$, apply $O*\,$ to get $X*Y=(X+Y)*O=-(X+Y)$. Then
\[
P\oplus Q=O_1*(P*Q)=-\bigl(O_1+(-(P+Q))\bigr),
\]
using the first identity once more. This proves the parenthesized formula with no associativity. Under the explicitly quoted abelian group theorem it simplifies to $P+Q-O_1$. The map from the new-origin group to the old one is $\phi(P)=P-O_1$, since
$\phi(P\oplus Q)=P+Q-2O_1=\phi(P)+\phi(Q)$. Its inverse is translation by $O_1$. This final group interpretation is conditional on that quoted theorem.
\end{solution}
\begin{solution}{geometry:ex:sixonconic}
\idea The conic and the line through two remaining intersections form a cubic through eight points.
\writeup The line-union cubic $P_1P_2\cup P_3P_4\cup P_5P_6$ has no common component with the smooth irreducible $E$, and its nine intersections are exactly the stated distinct points. The cubic $C\cup AB$ contains the six points on $C$ and also $A,B$, so Cayley--Bacharach puts $C'$ on it. Since $C'\notin C$, it lies on $AB$. The earlier collinearity exercise gives $(A+B)+C'=O$ directly. Also the three paired sums are $-A,-B,-C'$. Turning these into an unparenthesized sum of all six points requires the full associativity theorem quoted from Volume~13; the argument here proves the stated parenthesized conclusion.
\end{solution}

\begin{solution}{ex:v9-hesse-pencil}
\idea The four reducible members display all twelve lines, while a tangent substitution shows that the common basepoints stay flexes as the curve varies.
\writeup The gradient equations, divided by three, are $X^2=\kappa YZ$, $Y^2=\kappa XZ$, $Z^2=\kappa XY$. A zero coordinate would force all coordinates zero. Otherwise multiplying gives $\kappa^3=1$. For each such parameter a change $Z_{\rm new}=\kappa Z$ reduces the equation to the factorization
\[
X^3+Y^3+Z^3-3XYZ
=(X+Y+Z)(X+\omega Y+\omega^2Z)(X+\omega^2Y+\omega Z).
\]
The three lines are distinct and their intersection vertices are singular, so precisely these parameters are singular. Writing the pencil homogeneously in its parameter gives the fourth triangle $XYZ=0$ at infinity.

The common coordinate-line points are $[1:-u:0]$ and their coordinate permutations, with $u^3=1$, nine points in all. At $P=[1:-u:0]$ the tangent for a finite smooth parameter is $X+u^2Y+\kappa uZ=0$. On its parametrization $[1:-u-\kappa u^2t:t]$, substitution gives $(1-\kappa^3)t^3$. Thus the tangent has exactly triple contact at $P$. Symmetry gives all nine flexes. One flex allows the normal-form theorem to apply, and the nine-flex theorem makes this list exhaustive.

Each side of a finite triangle has one basepoint on each coordinate line, hence three. The infinity triangle's sides have their three coordinate-line basepoints. All twelve lines are distinct: a line shared by two different finite members would divide their difference, a nonzero multiple of $XYZ$, but no finite member has a coordinate-line component. Different lines share at most one basepoint. Their $12\binom32=36=\binom92$ pairs therefore account for every pair exactly once. In one triangle its three sides partition the basepoints: their vertices are not basepoints, as the displayed gradient at each basepoint is nonzero even for a singular parameter. A side of another triangle meets each of those three parts once, since its three basepoints cannot put two on the same side. These four parallel classes of three lines give exactly the affine-plane incidence pattern of order three. The proof uses equations and counting throughout.
\end{solution}
\begin{solution}{ex:v9-quadratic-twist-count}
\idea A nonsquare multiplier reverses every nonzero quadratic-residue sign.
\writeup For each $x$, the two ordinate counts are $1+\chi(f(x))$ and $1+\chi(f(x)/d)=1-\chi(f(x))$, with $\chi(0)=0$. Their sum is two. Both projective curves have exactly one infinity point, so the total sum of point counts is $2p+2$. Smoothness follows from squarefreeness just as for the original curve.

For the given example the values of $f$ are $1,3,1,1,4$. Dividing by two gives $3,4,3,3,2$. Only the value four is a square, at $x=1$, so the twist points are $(1,2),(1,-2),O$, three in all, and $9+3=12=2(5+1)$. Over the quadratic extension choose $c$ with $c^2=d$. The map $[x:y:z]\mapsto[x:cy:z]$ takes $dy^2z=f(x,z)$ to $y^2z=f(x,z)$ and has inverse scaling by $c^{-1}$. It is defined over that extension and fixes $O$. An isomorphism over $\mathbb F_5$ would biject rational points, impossible for counts three and nine. Thus extending the field matters.
\end{solution}
\begin{solution}{ex:v9-denominator-doubling}
\idea A uniquely least power of two cannot cancel, so the doubling formula forces ever larger denominators.
\writeup The tangent slope is $(3x^2+a)/(2y)$, giving $x_D=(3x^2+a)^2/(4f(x))-2x$. Combining terms yields the stated numerator $x^4-2ax^2-8bx+a^2$. Put $v=v_2(x)<0$, ignoring any zero terms in valuation comparisons. Because $a,b$ are integers, the term $x^3$ has uniquely least valuation in $f(x)$: $3v<v+v_2(a)$ and $3v<v_2(b)$ whenever those terms are present. Hence $2v_2(y)=3v$, so $v$ is even and $v_2(y)=3v/2$; in particular $y\ne0$.

In the numerator the valuation $4v$ of $x^4$ is strictly below the possible values $1+v_2(a)+2v$, $3+v_2(b)+v$, $2v_2(a)$ of the other terms. Thus its valuation is $4v$, while the denominator's is $2+3v$. Their difference is $v-2$. The output is another rational curve point by the chord construction, with negative abscissa valuation again. Induction therefore gives defined, distinct iterates, since their valuations are $v-2n$ and their ordinate valuations $3(v-2n)/2$. No reassociation of sums was used.

For the starting double of Bachet's point, $x=129/100$ has valuation $-2$ and $y=-383/1000$ valuation $-3$. After $n$ further doublings these become $-2-2n$ and $-3-3n$. This proves infinitude directly from elementary divisibility.
\end{solution}
\begin{solution}{ex:v9-nodal-multiplication}
\idea The third-intersection rule becomes multiplication after a fractional linear change of parameter.
\writeup The affine node is $(0,0)$. Outside it, $x\ne0$, and $t=y/x$ gives $x=t^2-1,y=t(t^2-1)$, with $t\ne\pm1$. The projective parameter infinity gives $O$, smooth with tangent $z=0$ and triple contact. The derivatives, or this inverse parameter at finite points, show these are all smooth points; the only excluded point is the node. The map $u=(t+1)/(t-1)$ identifies $\mathbb P^1(K)\setminus\{1,-1\}$ with $K^\times$, taking infinity to one. Its inverse is $t=(u+1)/(u-1)$, with $u=1$ giving infinity.

Every line through the node has intersection order at least two there, since the affine equation has no linear term. It cannot also have two smooth intersections counted with multiplicity. Thus a chord or tangent determined by two smooth points avoids the node, and its third point is smooth. For a nonvertical line $y=mx+b$, substitution of the parameter gives
\[
h(t)=t^3-mt^2-t+m-b.
\]
Such a chord cannot pass through the node, so $b\ne0$, and its three parameter roots $t_i$, with multiplicity, satisfy $h(1)=h(-1)=-b$. Therefore
\[
\prod_i\frac{t_i+1}{t_i-1}=\frac{h(-1)}{h(1)}=1.
\]
The signs from the three factors cancel between numerator and denominator. This identity includes tangent repetitions. For a vertical chord, the three parameters are $t,-t,\infty$, with product $u(t)u(-t)u(\infty)=1$; at $t=0$ it is the double finite contact and $O$. The tangent at $O$ has three copies of infinity and gives the same product one. These cases cover all lines used by the operation.

Consequently $u(P*Q)=(u(P)u(Q))^{-1}$. Since $u(O)=1$, applying $O*\,$ inverts again, giving $u(P+Q)=u(P)u(Q)$. This bijection transports every group axiom from $K^\times$ and proves associativity on the smooth locus directly. The singular point was excluded for a mathematical reason: lines through it spend two of the three intersections there and do not support this same smooth-point rule.
\end{solution}

\section*{Sources and further reading}

Silverman and Tate develop rational points and the full group law from these constructions. For the Hesse pencil, Artebani and Dolgachev give both the equations and the configuration's history. Kirwan and Brieskorn--Kn\"orrer explain Hessians and complex flexes. The original records below support the historical statements and identify the deeper arithmetic theorems quoted here.

\begin{itemize}
\item D.~Perrin, \emph{Autour de l'\'equation de Bachet}, lecture notes, Universit\'e Paris-Saclay (Bachet's tangent method of 1621 and the point $\bigl(\frac{129}{100},\frac{383}{1000}\bigr)$
of $y^2+2=x^3$).
\item I.~Newton, \emph{Enumeratio linearum tertii ordinis}, appendix to \emph{Opticks}, London, 1704; the description of Newton's manuscripts at the
Cambridge Digital Library (72 species, projections of the five divergent parabolas).
\item M.~Artebani and I.~Dolgachev, \emph{The Hesse pencil of plane cubic curves}, L'Enseignement Math\'ematique 55 (2009), 235--273 (the Hesse
configuration, its history and its symmetries).
\item J.~H.~Silverman and J.~Tate, \emph{Rational Points on Elliptic Curves}, 2nd ed., Springer, 2015 (the group law, Nagell--Lutz, Mordell's theorem and
curves over finite fields, at the level of this book).
\item H.~Poincar\'e, \emph{Sur les propri\'et\'es arithm\'etiques des courbes alg\'ebriques}, J.\ Math.\ Pures Appl.\ (5) 7 (1901); L.~J.~Mordell, \emph{On the
rational solutions of the indeterminate equations of the third and fourth degrees}, Proc.\ Cambridge Philos.\ Soc.\ 21 (1922); B.~Mazur, \emph{Modular
curves and the Eisenstein ideal}, Publ.\ Math.\ IH\'ES 47 (1977).
\item N.~Koblitz, \emph{Elliptic curve cryptosystems}, Math.\ Comp.\ 48 (1987); V.~Miller, \emph{Use of elliptic curves in cryptography}, CRYPTO '85 (1986).
\item E.~Brieskorn and H.~Kn\"orrer, \emph{Plane Algebraic Curves}, Birkh\"auser, 1986; F.~Kirwan, \emph{Complex Algebraic Curves}, Cambridge University Press,
1992 (the Hessian and the nine flexes).
\end{itemize}
